\documentclass[12pt, a4paper]{academics}

\academicsCourse{数据库原理与应用}{Database Principle and Application}
\academicsReport{需求分析——招投标文件}{2026 年 3 月 25 日}
\academicsArthur{华博文}{19240212}

\begin{document}

\makeacademicscover

\section{实验环境}

\subsection{操作系统}

macOS 26.4 arm64，Darwin 25.4.0，Apple M5。

\subsection{协作工具}

DeepSeek（业务顾问，扮演甲方时使用），豆包（技术顾问，扮演乙方时使用）。

\subsection{文档工具}

\LaTeX{}。

\section{实验目的}

本次实验旨在通过模拟真实软件开发项目的需求分析全流程，系统性地掌握数据库应用系统开发的方法论基础。实验的核心目标包含三个层面。首先是认知层面——理解需求分析在数据库应用系统开发项目中的关键地位：它是后续所有设计、实现与测试工作的出发点，需求阶段的任何偏差都将在系统上线后被成倍放大，纠正成本呈指数级增长。其次是方法层面——掌握通过角色扮演与大型语言模型协作进行业务需求梳理的方法，并能够将模糊的业务诉求转化为结构化的技术服务需求，这一能力在实际项目的需求沟通与方案汇报环节中至关重要。最后是实践层面——能够独立完成项目招投标过程中的核心文档撰写，包括业务需求文档、方案建议书以及数据字典，形成从需求分析到方案输出的完整闭环。

\section{实验内容}

本次实验选择「教务管理系统」作为项目主题，采用角色扮演的方式，分别模拟甲方（业务方）与乙方（技术方）两个立场，与不同的大型语言模型展开协作，完整经历从业务需求提出到技术方案设计的全过程。甲方角色使用 DeepSeek 作为业务顾问，乙方角色使用豆包作为技术顾问，两方各进行三轮迭代对话，共计六轮。每轮对话均保存了完整的提示词、LLM 输出原文及对话截图，以忠实记录需求逐步收敛、方案反复打磨的迭代轨迹。

\subsection{阶段一：扮演甲方（业务方）}

在阶段一中，我以高校教务管理部门的身份，借助 DeepSeek 进行业务需求梳理。三轮对话遵循"框架搭建—逐项细化—交付成型"的递进逻辑，逐步将模糊的建设愿景转化为结构完整、可直接用于招标的业务需求文档。

\subsubsection{Round 1：项目启动与需求框架}

首轮对话以宏观视角切入。我向 DeepSeek 描述了高校教务管理的典型场景，要求其以信息化咨询顾问的身份，为教务管理系统建设项目梳理项目背景、建设目标与核心功能模块。这一轮提示词的设计刻意保持开放，不限定具体格式与深度，目的在于观察 LLM 在无约束条件下能够建立多大规模的问题框架，同时检验模型在缺乏明确指令时自主进行结构化输出的能力。

\begin{center}
\includegraphics[width=\textwidth]{../res/phase1-round1.png}
\end{center}

首轮输出即产出了一份结构完整、逻辑清晰的《需求建议书》框架，超出预期。DeepSeek 通过对高校教务管理现状的系统分析，精准识别了六大核心痛点：系统架构陈旧难以扩展、各业务系统间数据孤岛严重、教学改革缺乏信息系统支撑、师生用户操作体验不佳、管理层缺乏数据驱动的决策工具、以及系统安全与合规方面存在明显风险。在此基础上，模型进一步提出了五条可量化、可验收的建设目标——支撑 30000 名学生并发选课、核心业务线上化率达到 100\%、跨部门数据共享延迟控制在 5 分钟以内、提供 20 类标准化报表、师生整体满意度超过 85\%——并初步规划了涵盖学籍管理、培养方案、排课管理、选课管理、成绩管理等在内的 12 个核心功能模块。这一轮输出的最大价值在于建立了一个共识框架，使得甲方与顾问在项目的边界与核心目标上迅速达成一致，为后续的深化细化奠定了坚实的讨论基础。

\subsubsection{Round 2：需求细化与非功能需求}

第二轮在首轮共识框架的基础上进行纵向深化。提示词要求 DeepSeek 对上一轮划分的 12 个功能模块逐项展开，每个模块至少给出 4 至 5 条可操作的具体功能描述，同时补充非功能性需求维度和面向投标方的资质要求。这一轮的核心策略是将"做什么"从概念层面落到执行层面，使得输出从"可以理解"走向"可以直接招标"。

\begin{center}
\includegraphics[width=\textwidth]{../res/phase1-round2.png}
\end{center}

本轮的输出显著提升了文档的可操作性与完整度。12 个功能模块均获得了详尽的子功能拆解——例如学籍管理模块细化为新生入学注册、学籍信息维护、学籍异动（休学、复学、转专业、退学）处理以及学籍状态查询与统计等具体场景，各子功能点均附有输入条件、处理规则和输出结果的描述，具备了直接转换为用例规格说明的基础。非功能性需求方面，模型从四个维度给出了可量化的约束指标：性能维度要求一般查询响应时间不超过 1 秒、95\% 的数据库事务在 5 秒内完成；安全维度明确需达到信息安全等级保护二级标准、全链路启用 HTTPS 加密传输；可用性维度提出了 99.95\% 的系统可用率目标，以及恢复点目标（RPO）小于 15 分钟、恢复时间目标（RTO）小于 1 小时的灾备指标；可扩展性维度则要求在架构层面预留与校园一卡通、图书馆管理系统、财务系统等既有系统的标准对接接口。此外，本轮还补全了投标方资质门槛、交付物清单、五阶段工期里程碑安排以及售后服务承诺等招标文件的结构性章节。至此，文档从首轮的"需求框架"提升到了"可发布招标"的成熟度。

\subsubsection{Round 3：业务需求文档}

第三轮进行了一次关键的视角转换。前两轮的输出本质上仍然是面向招标场景的需求建议书，而第三轮的目标是产出一份面向开发团队的业务需求文档——这份文档需要回答的核心问题是"系统应该为谁解决什么问题"，而非"应该采购怎样的系统"。提示词要求 DeepSeek 从最终用户的角度出发，系统性地梳理用户角色、核心业务流程、业务规则以及预期的业务价值。

\begin{center}
\includegraphics[width=\textwidth]{../res/phase1-round3.png}
\end{center}

DeepSeek 对本轮提示词的理解十分到位。在用户角色分析方面，它定义了覆盖教务管理全场景的 8 类角色——从最直接的学生和教师，到教务管理员、辅导员、教学院长与校领导等管理层，再到系统管理员和外部系统接口用户等支撑性角色——每一类角色均附有明确的使用场景描述和权限边界界定，为后续的权限模型设计提供了清晰的输入。在业务流程方面，模型以 Mermaid 流程图配合结构化文字，详述了学生选课、教师成绩录入、学籍异动这三个最核心也最复杂的端到端流程，涵盖正常路径与异常路径的完整处理逻辑，流程图的引入使得业务方与开发方可以在同一张图上对齐理解。尤为值得关注的是模型提炼出的 6 条不可绕过的业务规则——培养方案刚性约束、多维度实时冲突检测、人—时—空三不冲突的排课原则、成绩提交即锁定且修改必留痕的审计要求、自动批量毕业审核规则、以及三权分立的权限管控模型——这些规则的提炼体现了 LLM 在特定领域内进行跨轮抽象归纳的能力，也是人工难以在短时间内完成的系统性工作。最后，模型还从提升管理效率、保障教学质量、改善用户体验、支撑科学决策、确保合规安全五个维度阐述了项目的预期业务价值，为这份需求文档赋予了战略层面的说服力。至此，甲方三轮输出形成了从"建议书"到"详细需求"再到"面向开发的业务文档"的完整需求链条。

\subsection{阶段二：扮演乙方（技术方）}

阶段二中，角色切换至技术方。我将甲方三轮输出的完整内容提炼为核心摘要（涵盖项目背景、12 个功能模块需求、非功能性指标及数据量预估），作为唯一输入提供给豆包，要求其以技术架构师的身份进行方案应标。这一阶段模拟了真实招投标过程中的信息不对称——乙方只能基于甲方公开发布的招标文件进行方案设计，无法获知甲方的内部讨论和隐性偏好，因此对需求摘要的准确提炼与有效传递本身就构成了实验的一项重要变量。豆包同样进行了三轮迭代，分别完成了总体技术方案设计、实施计划与数据字典编制、以及安全策略与测试方案的查漏补缺。

\subsubsection{Round 1：技术方案设计}

首轮要求豆包基于甲方需求摘要，输出一份完整的技术方案建议书，涵盖总体架构设计、技术选型依据、12 个模块的核心子模块设计以及关键技术难点的应对策略。

\begin{center}
\includegraphics[width=\textwidth]{../res/phase2-round1.png}
\end{center}

豆包给出的技术方案在架构选型和业务匹配度上均表现出了较强的专业性。在总体架构层面，模型推荐采用微服务架构配合前后端分离的 B/S 模式，这一选择与教务管理系统多模块独立部署、按需扩展的天然需求高度吻合。前端技术栈选取了 Vue 3 + TypeScript + Element Plus 的组合用于 PC 端管理后台，并通过 uni-app 框架实现移动端覆盖，兼顾了开发效率与跨平台需求。后端基于 Spring Cloud Alibaba 微服务生态构建，以 Nacos 作为注册与配置中心、Gateway 统一网关入口、Sentinel 负责流量控制与熔断降级，形成了一套完整的服务治理体系。在数据存储层面，方案采用了多引擎组合策略：以 MySQL 8.0 作为核心业务数据库承担在线事务处理负载，Redis 7.x 构建分布式缓存集群以应对高并发读取场景，ClickHouse 负责统计分析类联机分析处理（OLAP）查询，Elasticsearch 则支撑全局搜索与日志检索需求——每种引擎各司其职，避免了单一数据库在混合负载下的性能瓶颈。

尤为突出的是选课高并发场景的应对设计。豆包提出了一套贯穿全链路的六层保障策略：网关层通过令牌桶算法实施限流削峰，防止瞬时流量击穿后端；Web 层设置本地缓存和 CDN 静态化以降低动态请求比例；服务层采用 Redis 缓存前置分担 90\% 以上的读压力，并通过分布式锁杜绝超选；消息队列（RocketMQ）实现异步落库，将同步的写峰值削平为异步的稳态流量；数据库层以乐观锁机制控制课程容量的最终一致性，确保不会出现一个名额被两人同时占据的情况；最后是降级熔断兜底——当任一环节压力超过预设阈值时自动触发服务降级，保证核心选课功能可用而非整个系统崩溃。这套策略从入口到存储形成了纵深防护，是本次实验中 LLM 在系统设计层面产出的高质量内容之一。对甲方的 12 个功能模块，模型逐一给出了核心子模块的技术设计方案，并以 Mermaid 依赖关系图直观展示了各模块间的数据流向与接口调用关系。

\subsubsection{Round 2：实施计划、部署方案与数据字典}

第二轮要求豆包从"方案"走向"落地"，补充项目的实施路线图、生产环境部署方案以及最核心的成果——数据字典定义。

\begin{center}
\includegraphics[width=\textwidth]{../res/phase2-round2.png}
\end{center}

在实施计划方面，豆包制定了一个为期 8 个月的六阶段推进方案，采用"里程碑式管控、分阶段验收"的管理模式。第 1 至 2 月为需求调研与系统设计阶段，交付物为需求规格说明书、架构设计说明书、数据库设计说明书及接口规范文档，须经甲方评审签字确认后方可进入下一阶段。第 3 至 5 月集中进行核心业务模块开发，以学籍管理、培养方案、排课、选课、成绩五个核心模块为先导，同步推进系统管理与权限框架的建设。第 6 月进入全模块集成与测试阶段，完成剩余七个模块的开发并进行系统性集成测试、性能压测和安全合规测试。第 7 月上半月执行生产环境部署与历史数据迁移，下半月启动为期一个月的试运行与校院两级用户培训。第 8 月下半月完成最终验收与运维交接。方案同时配置了涵盖项目经理、技术架构师、前后端开发工程师、测试工程师、实施顾问、UI/UX 设计师、运维工程师以及教务业务专家在内的 17 人项目团队，各角色职责清晰、覆盖从技术实现到业务咨询的全链路。

部署方案的设计充分体现了高可用与安全合规的工程考量。整个生产环境规划了约 40 个物理或虚拟节点，纵向划分为接入安全、应用服务、中间件、数据存储和运维监控五个层次，每一层均采用集群或主备模式部署以消除单点故障。横向则按照安全等级划分为五个网络隔离区域：外网接入区（防火墙、WAF、负载均衡）、DMZ 隔离区（门户网关、认证服务、API 网关）、业务服务区（全部微服务应用节点）、数据存储区（MySQL 主从集群、Redis 缓存集群、ClickHouse 分析集群、Elasticsearch 搜索集群、MinIO 对象存储）以及运维管理区（堡垒机、监控平台、日志审计系统），数据存储区仅允许来自业务服务区指定 IP 的访问，安全边界清晰且符合等保二级的纵深防护要求。

数据字典是这一轮最重要的产出。豆包给出了院系信息表、专业信息表、学生信息表、教师信息表、课程信息表、选课记录表、成绩信息表共 7 张核心业务表的完整字段定义，并确立了全系统的命名规范——统一以 \mintinline{text}{edu_} 和 \mintinline{text}{sys_} 前缀区分业务表与系统表，所有业务表均包含 \mintinline{text}{create_time}、\mintinline{text}{update_time}、\mintinline{text}{is_deleted} 三个通用字段，分别用于记录创建时间、最后修改时间以及逻辑删除标记。这一命名与字段规范看似简单，但在实际项目中对于降低团队沟通成本、保证数据库结构的长期可维护性以及支持审计追踪需求具有重要的工程价值。

以下为数据字典中 5 张核心业务表的字段定义：

\setlength{\belowcaptionskip}{8pt}
\begin{table}[h]
\centering
\caption{学生信息表（edu\_student）}
\begin{tabular}{|c|c|c|c|c|}
\hline
\textbf{字段名} & \textbf{数据类型} & \textbf{长度} & \textbf{键} & \textbf{说明} \\
\hline
student\_id & VARCHAR & 20 & 主键 & 学号 \\
student\_name & VARCHAR & 50 & -- & 学生姓名 \\
gender & CHAR & 1 & -- & 性别（M/F） \\
id\_card & VARCHAR & 18 & -- & 身份证号 \\
department\_id & VARCHAR & 32 & 外键 & 所属院系ID \\
major\_id & VARCHAR & 32 & 外键 & 所属专业ID \\
grade & VARCHAR & 4 & -- & 入学年级 \\
student\_status & TINYINT & -- & -- & 学籍状态 \\
\hline
\end{tabular}
\end{table}

\begin{table}[h]
\centering
\caption{教师信息表（edu\_teacher）}
\begin{tabular}{|c|c|c|c|c|}
\hline
\textbf{字段名} & \textbf{数据类型} & \textbf{长度} & \textbf{键} & \textbf{说明} \\
\hline
teacher\_id & VARCHAR & 20 & 主键 & 教师工号 \\
teacher\_name & VARCHAR & 50 & -- & 教师姓名 \\
gender & CHAR & 1 & -- & 性别（M/F） \\
department\_id & VARCHAR & 32 & 外键 & 所属院系ID \\
title & VARCHAR & 30 & -- & 职称 \\
teacher\_status & TINYINT & -- & -- & 在职状态 \\
phone & VARCHAR & 20 & -- & 联系电话 \\
email & VARCHAR & 100 & -- & 电子邮箱 \\
\hline
\end{tabular}
\end{table}

\begin{table}[h]
\centering
\caption{课程信息表（edu\_course）}
\begin{tabular}{|c|c|c|c|c|}
\hline
\textbf{字段名} & \textbf{数据类型} & \textbf{长度} & \textbf{键} & \textbf{说明} \\
\hline
course\_id & VARCHAR & 32 & 主键 & 课程唯一ID \\
course\_code & VARCHAR & 20 & -- & 课程编码 \\
course\_name & VARCHAR & 100 & -- & 课程名称 \\
credit & DECIMAL(3,1) & -- & -- & 学分 \\
total\_hours & INT & -- & -- & 总学时 \\
course\_type & VARCHAR & 20 & -- & 课程类型 \\
department\_id & VARCHAR & 32 & 外键 & 开课院系ID \\
prerequisite & VARCHAR & 255 & -- & 先修课程要求 \\
\hline
\end{tabular}
\end{table}

\begin{table}[h]
\centering
\caption{选课记录表（edu\_enrollment）}
\begin{tabular}{|c|c|c|c|c|}
\hline
\textbf{字段名} & \textbf{数据类型} & \textbf{长度} & \textbf{键} & \textbf{说明} \\
\hline
enrollment\_id & BIGINT & -- & 主键 & 雪花算法主键 \\
student\_id & VARCHAR & 20 & 外键 & 学号 \\
course\_id & VARCHAR & 32 & 外键 & 课程ID \\
semester & VARCHAR & 20 & -- & 学年学期 \\
teacher\_id & VARCHAR & 20 & 外键 & 授课教师工号 \\
enroll\_type & TINYINT & -- & -- & 选课类型 \\
enroll\_status & TINYINT & -- & -- & 选课状态 \\
\hline
\end{tabular}
\end{table}

\begin{table}[h]
\centering
\caption{成绩信息表（edu\_score）}
\begin{tabular}{|c|c|c|c|c|}
\hline
\textbf{字段名} & \textbf{数据类型} & \textbf{长度} & \textbf{键} & \textbf{说明} \\
\hline
score\_id & BIGINT & -- & 主键 & 雪花算法主键 \\
enrollment\_id & BIGINT & -- & 外键 & 选课记录ID \\
student\_id & VARCHAR & 20 & 外键 & 学号 \\
course\_id & VARCHAR & 32 & 外键 & 课程ID \\
semester & VARCHAR & 20 & -- & 学年学期 \\
usual\_score & DECIMAL(5,2) & -- & -- & 平时成绩 \\
final\_score & DECIMAL(5,2) & -- & -- & 期末成绩 \\
total\_score & DECIMAL(5,2) & -- & -- & 总评成绩 \\
grade\_point & DECIMAL(3,2) & -- & -- & 对应绩点 \\
score\_status & TINYINT & -- & -- & 成绩状态 \\
\hline
\end{tabular}
\end{table}

\subsubsection{Round 3：数据字典完善与技术方案补充}

第三轮采取"自查补漏"的策略——要求豆包以审阅者的视角审视前两轮输出的完整性与一致性，主动识别缺失或不完备的内容，并从数据字典、安全策略、测试方案和风险分析四个维度进行系统性补充。这种让 LLM 自我审查的策略在本实验中表现出显著效果：模型在"生成"模式下倾向于聚焦主线而忽略边缘场景，而在"审阅"模式下则能有效地发现盲区。

\begin{center}
\includegraphics[width=\textwidth]{../res/phase2-round3.png}
\end{center}

豆包在本轮中展现出了较强的自查意识与补充能力。数据字典方面，它不仅给出了以 Mermaid 描述的实体关系模型（清晰展示 14 张表之间的外键依赖与一对多关联），还识别出前两轮遗漏的 7 张支撑表——学期表、系统用户表、角色表、用户角色关联表、菜单权限表、角色菜单关联表以及操作日志表——使数据库设计从"核心业务"层扩展至"系统支撑"层，形成了从业务操作到权限管控再到审计追踪的完整数据闭环。同时补充了 16 个核心索引的定义：唯一索引覆盖院系编码、专业编码、身份证号、课程编码等需要全局唯一性保证的字段；组合索引针对高频查询场景（如按院系—专业—年级联合查询学生、按学生—学期联合查询选课记录、按课程—学期联合查询成绩）进行了优化；审计类索引则为操作日志表建立了操作人—时间的复合索引，以支持高效的审计追查。索引设计的合理性体现了模型对实际查询负载特征的理解——并非简单地给每个字段建索引，而是在冗余与查询效率之间做出了取舍。

安全策略部分，豆包给出了层次分明的纵深防护方案。在应用层，通过 SQL 注入三级防护体系——代码层 MyBatis 参数预编译（杜绝字符串拼接）、网关层 WAF 规则库实时拦截（目标准确率不低于 99.9\%）、数据库层最小权限原则与微服务独立账号——形成多重阻断。敏感数据实施全生命周期保护：身份证号与手机号采用 AES-256 加密存储，密码使用 BCrypt 哈希加盐，传输强制使用全链路 HTTPS/TLS 1.3，展示侧根据角色实施分级脱敏，审计侧记录所有敏感数据的全量访问日志。备份体系按照恢复时间与留存需求分层设计：MySQL 核心数据以主从半同步复制实现实时冗余，每日全量备份与每小时增量日志备份构成 30 天的可追溯窗口，每月冷备归档满足 3 年的长期合规留存要求。

测试方案规划了五个递进阶段：JUnit 5 与 Mockito 驱动的单元测试确保核心业务逻辑覆盖率不低于 95\%；Postman 与 JMeter 执行的集成测试保证全部 API 接口通过率 100\%；JMeter 与 Gatling 实施的性能压测要求峰值 QPS 不低于 5500、错误率低于 0.1\% 且 72 小时连续运行无宕机；安全合规测试通过自动化扫描与人工渗透相结合确保高危和中危漏洞 100\% 清零并满足等保二级达标要求；最终的用户验收测试由甲方验收小组在真实业务场景下确保核心流程 100\% 跑通。风险分析部分，模型识别了选课高并发性能不达标、历史数据迁移失败、需求频繁变更、上线窗口与教学周期冲突、信创环境兼容性以及用户操作适配度低等六项核心风险，对每项风险给出了等级评定与可操作的应对措施，如针对选课高并发风险提出了三轮全链路压测、1.5 倍冗余配置和分级降级预案的组合应对策略。这一轮输出将整个技术方案从"功能完备"提升到了"可交付工程"的成熟度。

\section{实验过程归纳总结}

\subsection{多轮迭代策略}

回顾整个实验过程，最核心的方法论收获在于掌握了一套与 LLM 协作进行需求分析的渐进式迭代策略。无论是甲方还是乙方角色，实验均严格遵循三轮递进的对话结构，而这一策略的有效性在两阶段实验中得到了反复且一致的验证。

第一轮为框架层：以开放性问题引导 LLM 建立对问题域的整体认知，不限定输出格式与深度。这一阶段产出的内容覆盖面广但深度尚浅，其价值在于快速建立讨论的边界与共同语境——用项目管理的语言来说，类似于需求启动阶段的"项目章程"，使所有参与方在范围和目标上达成初步一致。第二轮为细化层：基于首轮确立的宏观框架，逐模块、逐维度地要求 LLM 展开与补充。甲方阶段将 12 个功能模块逐一拆解为可操作的功能点，并补全了性能、安全、可用性、可扩展性四个维度的非功能性约束；乙方阶段则将技术架构落地为可执行的项目时间表、可采购的硬件配置清单和可直接建表的数据字典。第三轮为交付层：通过主动的视角切换——甲方从招标方视角转为开发团队视角，乙方从方案设计者视角转为自我审阅者视角——促使 LLM 识别遗漏、补充边界条件，最终产出可直接交付使用的完整文档。

这一迭代策略之所以有效，根本原因在于 LLM 的注意力机制存在天然的"广度—深度"权衡。一次性要求输出全部内容时，模型倾向于给出面面俱到但每一点都浅尝辄止的泛化结果；而通过多轮对话将问题空间逐层收窄，每一轮都可以在前一轮成果的基础上进行更深层次、更聚焦的推理与展开。实验中的一个清晰对比是：如果直接要求 LLM"输出教务管理系统的完整需求文档"，得到的往往是一份正确但缺乏个性的模板化产物；而通过三轮递进，最终产出的文档不仅结构完整，还包含了从高校实际痛点推导出的业务规则、从并发场景反推出的技术策略，以及从自我审阅中发现并填补的数据模型盲区——这些"深一度"的内容正是迭代策略的核心价值所在。

\subsection{角色分离与上下文传递}

实验中刻意使用了两款不同的 LLM 分别扮演甲方（DeepSeek）和乙方（豆包），而非使用同一个模型完成全部对话。这一设计背后有两层考量，且均得到了实验结果的印证。

第一层考量是模拟真实招投标场景中的信息不对称。在真实项目中，乙方技术团队无法直接获知甲方的全部内部讨论、隐性需求与优先级偏好，只能基于甲方公开发布的招标文件进行方案设计。因此，实验中我需要将甲方三轮对话的完整输出提炼为一份结构化的需求摘要，并将其作为传递给乙方 LLM 的唯一上下文——这一"信息传递"环节本身即构成了实验的重要变量。提炼摘要的过程让我深刻体会到需求传递链条上信息压缩与信息保真之间永恒的矛盾：摘要过于简略会导致乙方缺乏足够上下文而产出泛化、缺乏针对性的方案；摘要过于冗长则失去了"提炼"的意义，且可能引入噪声干扰模型的判断。经过权衡，最终确定以"项目背景与痛点 + 12 个模块功能清单 + 四维非功能性指标 + 数据量级预估"四要素构成的摘要格式，这一粒度为两个极端之间找到了一个可用的平衡点。

第二层考量在于避免单一 LLM 可能产生的思维惯性与路径依赖。不同模型在训练数据分布、对齐策略和输出风格上的差异，在特定领域任务中可能转化为互补优势，而非简单的优劣之分。实验中观察到的现象印证了这一判断：DeepSeek 在结构化长文档的生成上表现更为突出，其输出的需求建议书和业务需求文档在章节组织、层次划分和逻辑递进方面均具备较高水准，且善于在文档中建立前后呼应的引用结构；而豆包在国内高校信息化技术栈的选型上展现出更强的场景适配能力——它推荐的 Spring Cloud Alibaba + Vue 3 + Element Plus 技术组合、40 节点五层网络拓扑部署方案以及等保二级合规策略，均体现了对国内高校实际 IT 环境与采购生态的深入理解。如果统一使用同一模型完成全部实验，这种角色适配带来的互补性将无从体现，乙方方案的"国内化"贴合度可能会大打折扣。

\subsection{人工判断的不可替代性}

尽管 LLM 在本实验中输出的内容在数量与质量上都达到了可观的水平——六轮对话累计产出了超过两万字的专业技术文档——但整个实验过程也明确地揭示出人工判断在三个关键环节上的不可替代性。

首先是需求合理性的校验环节。LLM 在生成量化指标时存在明显的"向上取整"倾向：例如首轮输出曾提出"支撑 10 万人并发选课"的性能目标，这对于一所普通规模的高校而言严重脱离实际。这并非模型的知识缺陷或幻觉，而是其训练语料中缺乏具体办学规模的信息，因此倾向于给出一个在所有可能情境下都"足够安全"的上界估计。实验中通过人工比对目标高校的实际在校生人数，将并发指标修正为 30000 人，这一干预看似微小，却是确保方案具备实际可操作性的前提——过高的性能要求将直接导致不必要的基础设施成本膨胀。

其次是技术选型的最终决策环节。LLM 善于罗列主流技术选项并分析各自的利弊权衡，但选型的最终判断需要综合考量项目预算约束、团队既有技术储备、与学校现有信息系统架构的兼容性、以及技术栈的长期维护可持续性等 LLM 无法获知的因素。实验中豆包推荐的技术栈整体合理且自成体系，但部分选型——例如 ClickHouse 分析型数据库的引入——在真实项目中需要结合学校的数据量级和分析需求的紧迫程度进行进一步论证，而非机械地采纳推荐。

最后是文档的结构化组织环节。LLM 的输出本质上是线性展开的文本流，而一份符合学术规范的实验报告需要在原始输出的基础上进行多维度的再加工：从海量信息中提取关键结论、按照报告的逻辑而非对话的顺序重新编排章节、在各轮输出之间建立横向对比与纵向因果关联、以及在每个小结中完成从"模型输出了什么事实"到"这些事实意味着什么"的分析跳跃。实验中六轮对话的原始输出总量庞大且存在内容重叠，如何从中甄别核心信息并组织成本报告的结构——这些工作均无法由 LLM 独立完成，而需要人的整体性判断和结构化编排能力。

\subsection{存在的问题与改进方向}

在肯定实验成果的同时，有必要诚实地审视其中暴露的问题与局限，这些反思对今后的类似实验具有更直接的参考价值。

其一，LLM 输出的一致性问题不容忽视。跨轮对话中偶有出现前后信息的矛盾：甲方 Round 1 与 Round 2 对个别模块的划分边界不完全一致，乙方 Round 1 与 Round 2 在部分技术参数上存在细微出入。这种不一致的根源在于当前主流 LLM 的无状态生成特性——每轮对话对于模型而言都是一次独立的生成任务，它并不真正"记忆"或"理解"上一轮的每一个细节，只能依赖上下文窗口中提供的文本进行推理，而上下文窗口越长，模型对其中细粒度信息的注意力就越稀释。实验中通过人工校核和对齐来处理这些差异，但在实际项目中，这一问题需要更系统性的应对策略，例如将关键决策点显式地提炼为"决策日志"并强制纳入每轮提示词，而非依赖模型自主维护上下文一致性。

其二，提示词的敏感度比预期更高。实验过程中观察到，提示词措辞的细微调整——例如将"请给出技术方案"改为"请以技术架构师的身份，面向高校信息化部门负责人，输出一份可提交评审的方案建议书"——可能导致输出在风格、篇幅甚至核心内容取舍上发生显著变化。这种高敏感度使得提示工程本身成为一项需要反复试错的经验性技能，而非一套可精确复制、可稳定迁移的标准操作流程。通过六轮实验的积累，确实逐步沉淀出了一些有效模式——如明确且具体的角色身份设定、提供输出结构蓝本而非仅描述期望、在追问中显式引用前轮的关键输出等——但这些经验的迁移有效性仍然具有较强的任务依赖性和模型依赖性。

其三，角色扮演方法存在固有的认知边界。LLM 对高校教务管理业务的所有"理解"均建立在训练语料的统计模式之上，而非真实世界中多年积累的领域经验与隐性知识。这意味着模型能够生成在形式上高度专业的文档——使用正确的术语、遵循标准的文档结构、给出合理的技术指标——但在涉及那些只有亲身经历过高校行政管理复杂性才能把握的环节时，模型的输出可能流于表面或过于理想化。这是当前 LLM 技术的一个基本限制，短期内难以通过提示工程或迭代策略完全突破。因此，在实际应用中，LLM 更适合被定位为需求分析的"加速器"和"结构化助手"——它能够极大地缩短从零到一份可评审初稿的时间，但这份初稿必须经过领域专家的实质性审阅和修正，才能最终形成可靠的工作产物。

\section{结论}

本实验通过角色扮演与大型语言模型协作的方式，完整模拟了数据库应用系统开发中需求分析阶段的核心工作流程，从甲方需求梳理到乙方方案应标，共完成六轮迭代对话，产出了包括需求建议书、业务需求文档、技术方案建议书、实施计划、部署方案、数据字典、安全策略、测试方案与风险分析在内的全套项目文档。

在扮演甲方角色时，借助 DeepSeek 进行了三轮递进式迭代：从识别痛点与建立目标的项目框架层出发，经过 12 个模块逐项展开与四维非功能性约束补全的细化层，最终抵达面向开发团队的业务需求文档交付层，形成了从宏观愿景到微观业务规则的完整需求链路。在扮演乙方角色时，将甲方输出提炼为结构化摘要并传递给豆包，同样经历三轮迭代，完成了微服务架构选型、40 节点五层网络拓扑集群部署设计、14 张核心数据表的字典定义，以及纵深安全防护策略、五阶段递进测试方案和六维系统化风险分析的全套技术方案编制。

回顾整个实验过程，最深切的体认可以凝练为一句话：需求分析从来不是一次性输出的产物，而是一个在反复对话与反馈中逐步收敛的过程；大型语言模型在这一过程中扮演着极为高效的协作工具角色，但方案的最终正确性与实际适用性，仍然取决于——且必须取决于——人的专业判断力与领域知识。通过本次实验掌握的「框架—细化—交付」三轮递进提示策略，以及「不同角色匹配不同模型以发挥互补优势」的协作模式，为今后从事实际数据库应用系统开发中的需求分析工作积累了切实可迁移的方法论基础。

\appendix
\catcode`\%=12
\section{附录：LLM 提示词与完整输出}

以下保存了实验中 6 轮对话的完整提示词与 LLM 输出结果。

\subsection{阶段一（甲方，DeepSeek）}

\subsubsection{Round 1：项目启动与需求框架}

\textbf{提示词：}
\inputminted[]{text}{../res/phase1-round1-prompt.md}

\textbf{输出：}
\inputminted[]{text}{../res/phase1-round1-output.md}

\subsubsection{Round 2：需求细化与非功能需求}

\textbf{提示词：}
\inputminted[]{text}{../res/phase1-round2-prompt.md}

\textbf{输出：}
\inputminted[]{text}{../res/phase1-round2-output.md}

\subsubsection{Round 3：业务需求文档}

\textbf{提示词：}
\inputminted[]{text}{../res/phase1-round3-prompt.md}

\textbf{输出：}
\inputminted[]{text}{../res/phase1-round3-output.md}

\subsection{阶段二（乙方，豆包）}

\subsubsection{Round 1：技术方案设计}

\textbf{提示词：}
\inputminted[]{text}{../res/phase2-round1-prompt.md}

\textbf{输出：}
\inputminted[]{text}{../res/phase2-round1-output.md}

\subsubsection{Round 2：实施计划、部署方案与数据字典}

\textbf{提示词：}
\inputminted[]{text}{../res/phase2-round2-prompt.md}

\textbf{输出：}
\inputminted[]{text}{../res/phase2-round2-output.md}

\subsubsection{Round 3：数据字典完善与技术方案补充}

\textbf{提示词：}
\inputminted[]{text}{../res/phase2-round3-prompt.md}

\textbf{输出：}
\inputminted[]{text}{../res/phase2-round3-output.md}

\end{document}
